\documentclass[11pt]{article}
\usepackage[margin=1in]{geometry}
\usepackage{amsmath,amssymb}
\title{P1 merge-scope note v1.3}
\date{19 September 2026}
\begin{document}
\maketitle
\section*{Purpose}
This note records the final scope corrections applied immediately before the merged manuscript.

\section*{Reliability-sign sectors}
For block reliabilities $r_\tau=1-2\eta_\tau$, the universal gauge acts by global negation
$r\mapsto-r$.  The relative signs
\[
\chi_\tau=\operatorname{sgn}(r_1r_\tau),\qquad \tau\ge2,
\]
are gauge-invariant and classify $2^{T-1}$ sign-pattern orbits.  The block-classification theorem is proved only in the same-sign orbit $\chi\equiv+1$, followed by the section $r_1>0$.  Mixed-sign sectors remain unclassified.  Thus, for $T\ge2$, the all-positive chamber is a substantive restriction plus gauge fixing; for $T=1$ it is pure gauge fixing.

\section*{Existential regularity}
The $(4)$ reconstruction does not require all $2|2$ flattenings to be nondegenerate.  It requires at least one valid grouping.  Likewise, the $(3,1)$ size-three reconstruction only requires at least one admissible pair choice with nonzero inversion denominator and a degree-one common-$\rho$ gcd.  The merged manuscript therefore treats these regularity conditions existentially.

\section*{Headline}
At $n=4$, in the same-sign sector and under partition-specific algebraic regularity, the fully singleton partition is generically locally identifiable with certified open-set global ambiguity, while every coarser partition is generically globally identifiable after choosing the positive-reliability gauge section.
\end{document}
